\scanpage{15}
Points fixes:
\[
H_0:\left\{\begin{aligned}(\lambda_{02}-1)x_i&=0\quad\text{si }2\leq i\leq n+1\\2x_1+(1-\lambda_{01})(x_2+\ldots+x_{n+1})&=1\end{aligned}\right.
\]
\deleted{\unsure{Si}} Si $\lambda_{0,2}-1$ inv. ces équations s'écrivent \deleted{$x_i=0\quad2\leq i\leq n+1$}
\[
\left\{\begin{aligned}x_i&=0\qquad2\leq i\leq n+1\\2x_1\deleted{\unsure{+}(1-\lambda_{01})(x_2+\ldots+x_{n+1})}&=1\end{aligned}\right.
\]
donc c'est \unsure{une|un} sous-espace linéaire affine de dim $0$ \deleted{\unsure{unique}} si $2$ inv. \deleted{(car. $\neq2$)}\inserted{\deleted{\unsure{si $k\ne0$|unique et non vide}} (car $\neq2$)}, et c'est vide si $2$ nul dans $k$ \inserted{[\deleted{\unsure{donc pas de}} \unsure{même} en car. $2$, l'hyperplan \unsure{à|:} l'infini\ldots)]} -- donc ce n'est pas un hyperplan. Donc il faut que $\lambda_{02}-1$ soit non inv. \inserted{i.e. nul} sur $\forall$ fibre, donc nilpotent, i.e. $\lambda_{02}=1+\unsure{\nu|v}$, $\nu$ nilpotent. Alors les équations deviennent
\[
\left\{\begin{aligned}\nu x_i&=0\qquad2\leq i\leq n+1\\2x_1+(1-\lambda_{01})(x_2+\ldots+x_{n+1})&=1\end{aligned}\right.
\]
\unsure{que c'est|qui est} un hyperplan sur \uline{chaque fibre}, sauf en \deleted{car} car. $2$ si $\lambda_{0,1}=1$, où c'est vide (mais alors on trouve l'hyperplan \unsure{à|:} l'infini).

\margin{introduire la variable d'homogénéité $\unsure{z|2}$ d'où équation $2x_1+(1-\lambda_{0,1})(x_2\ldots+x_{n+1})=z$}
% slanted

Pour que ce soit un hyperplan, \deleted{il f. et s.} \deleted{\unsure{on} si $2$ inv.} aux pts où la car. est $\neq2$ \deleted{\unsure{et bien}} ou $\lambda_{01}-1$ inv., comme \deleted{l'hyperplan est} $H_0$ contenu dans $H'_0$ d'équation $2x_1+(1-\lambda_{01})(x_2\ldots x_{n+1})=z$, il faut qu'on ait égalité; donc que les équations $\nu x_i=0$ \inserted{$(2\leq i\leq n+1)$ (après $\forall$ chgt. de base\ldots)} soient conséquence de la précédente.

\deleted{Si $2$ inv. ceci} \inserted{\unsure{Ceci exige|Cela exige}} exige \inserted{évidemment} que $\nu=0$.

\deleted{Si $2$ non inv. mais $(1-\lambda_{01})$ inv., alors pour $\uline{n\geq2}$ on trouve ces \unsure{conclusions|conditions}. Le cas $n=1$ doit être traité à part\ldots Si $2$ \uline{\unsure{et|est}} $(1-\lambda_{0,\unsure{1|i}})$}
% archive sheet number: 14
